\section{Results}
\label{sec:results}

\subsection{The body: growth, persistence, regeneration}
\label{sec:results-morph}

The substrate must first do the known thing before it can be asked to do the new
one. A single organism was trained for 6000 steps (94 minutes on an Apple M2,
seed 0, no death-guard resets), growing for the first 1000 steps and then
training persistence and regeneration from a sample pool.

Both morphology gates pass (Table~\ref{tab:g1g2}). Growth reaches
$\IoU = 0.9898$ at the training horizon $t=64$, with a 5th percentile of
$0.9826$ over 128 rollouts, so the result is not carried by a good average over
bad tails.

The more informative number is $\IoU = 0.9993$ at $t=256$. The organism is
\emph{better} at four times the training horizon than at the horizon itself, and
holds its area to within 0.7\% of target. The target morphology is therefore a
genuine attractor of the learned dynamics rather than a memorised trajectory that
decays once the training window ends --- which is the property everything
downstream depends on, since the damage experiment reads out long after $T_{\mathrm{grow}}$.

Regeneration is unambiguous: after destroying a disk containing 30\% of alive
cells, \textbf{all 256 trials recovered}, under the strict criterion of sustained
recovery over 20 consecutive steps, with a median recovery time of 13 steps.
Recovery was equally complete whether or not the destroyed disk covered the grid
centre, which is worth recording ahead of the sensor-targeted damage arm: at this
stage, position of damage does not matter for the body.

\begin{table}[t]
\centering
\small
\caption{Morphology gates. All values measured on the committed checkpoint;
G1 over 128 rollouts, G2 over 256 damage trials.}
\label{tab:g1g2}
\begin{tabular}{@{}llrl@{}}
\toprule
& \textbf{Metric} & \textbf{Value} & \textbf{Threshold} \\
\midrule
\multirow{5}{*}{G1}
& IoU at $t=64$ (mean)          & 0.9898 & $\geq 0.90$ \\
& IoU at $t=64$ (5th pct.)      & 0.9826 & $\geq 0.85$ \\
& RMSE at $t=64$                & 0.0109 & $\leq 0.05$ \\
& IoU at $t=256$                & 0.9993 & $\geq 0.85$ \\
& area ratio $|A_{256}|/|A_{64}|$ & 0.9931 & $[0.9, 1.1]$ \\
\midrule
\multirow{4}{*}{G2}
& fraction recovered            & 1.0000 & $\geq 0.90$ \\
& sustained 20 steps            & 1.0000 & reported \\
& median recovery time          & 13 steps & --- \\
& achieved severity             & 0.2982 & $0.30 \pm 0.02$ \\
\bottomrule
\end{tabular}
\end{table}

\subsection{What the substrate transports on its own}
\label{sec:results-probe}

Everything in Section~\ref{sec:task} rests on an assumption that had not been
tested: that a signal injected at the sensor patch can cross the body at all. We
measured it directly on the morphology-trained organism above, whose maximum
geodesic radius is 9 cells.

The probe injects each of the 8 bipolar referent codes into the sensor patch and
measures, per cell and per step, the fraction of state variance explained by
which code was injected. Bodies are grown once and shared across codes so the
morphology is bit-identical, and the stochastic fire masks are shared across
codes as well --- without that, between-code variance is contaminated by update
noise rather than reflecting the injected signal. Distance is Chebyshev geodesic
distance on the alive mask, which is exactly the NCA light cone.

The result is stark. Variance explained collapses with distance:
$0.951$ at the sensor patch, $0.494$ at one cell, $0.174$ at two, $0.039$ at
three, and indistinguishable from zero by five. The propagation front, defined as
the furthest ring exceeding a 5\% threshold, \textbf{stalls at $d=2$ from
$t \approx 24$ onward and never advances again}, at a fitted speed of $0.039$
cells per step against a body of radius 9. The body is never covered.

\gated{the ridge-decoder analysis reported in the milestone notes ---
per-ring decoding accuracy showing information reaching $d \approx 5$ rather than
$d \approx 3$ --- has no committed implementation and cannot be regenerated from
the saved probe artifact, which stores variance ratios rather than per-cell
states. It is excluded from this paper until the decoder is written into the
probe and re-run. Do not restore the table without the code}

\paragraph{Why this is a lower bound, and what follows.} This organism was
trained only to hold a shape. Nothing in its objective ever referenced the sensor
channels, so it had no reason to relay anything; the measured range is what the
substrate does \emph{incidentally}. The honest reading is that passive dynamics
carry a referent partway across the body and the task must supply the rest.

That reading changed three design decisions. It makes the per-cell vote loss
load-bearing rather than merely an anti-centralisation device --- supervising
distant cells is the only pressure that demands propagation, since a distant cell
can be correct only if the signal arrived. It sets the body radius to 7 rather
than 10, inside the measured envelope. And it motivates the broadcast precursor
below, which tests the propagation question before paying for the full game.

Building the probe before the game is what surfaced this. Measured instead as
accuracy against rollout length, the same failure would have looked like ``needs
more steps,'' and the fix would have been to lengthen the rollout --- against a
problem that time alone does not solve.

\subsection{Broadcast: buying propagation with a per-cell loss}
\label{sec:results-broadcast}

\gated{this run is still in progress at the time of writing; numbers below are
from a live training log and will be replaced with a completed run and a proper
held-out evaluation}

The broadcast task supplies exactly the pressure the probe showed was missing.
Warm-started from the morphology checkpoint above, with a body of radius 7 and
every alive cell supervised to vote for the injected referent, per-cell vote
accuracy rises from chance ($0.125$ on 8 referents) to above $0.90$, with
cell-level consensus rising in step from $\approx0.25$ to $\approx0.89$ while the
body holds at $\IoU \approx 0.88$.

Consensus and accuracy rising together is the meaningful part. A megaphone
solution would show accuracy rising while consensus stayed low, since a small
loud region can carry the pooled decision without the rest of the body knowing
anything. \gated{quorum fraction is implemented but has never been measured on a
trained organism; it is the direct test of that claim and must be reported here}

\paragraph{A measured conflict between body and message.} The weighting between
the morphology loss and the vote loss is not a free parameter. At equal weight,
the vote gradient dominates and the organism dissolves itself: accuracy climbs
while $\IoU$ falls by half, trading the body for the message. The loss
\emph{values} understate this badly --- they differ by an order of magnitude in
the opposite direction --- so the balance has to be set by watching per-term
gradient norms, not losses. \gated{this observation comes from a run that was
deleted rather than archived, and the gradient-ratio logger has an off-by-one
that records NaN on every logged step; both must be fixed and the sweep re-run
before this paragraph can carry a number}

\subsection{Protocol emergence}
\label{sec:results-emergence}

\gated{requires the full Lewis game with a receiver, a discrete channel, and
cross-play across seeds. Gates G3 and G4 are specified in
Section~\ref{sec:protocol} but not yet implemented. This section reports
accuracy with Wilson intervals, debiased $\MI(M;R)$, effective vocabulary, the
four mandatory interventions, causal influence of communication, and the
emergence rate across seeds}

\subsection{The dissociation}
\label{sec:results-dissociation}

\gated{the central experiment. Damage the sender with the receiver frozen and
healthy, and measure morphological against semantic recovery on the same axis.
Reports the recovery time courses, the dose-response map with the dissociation
quadrant of Eq.~\ref{eq:h3} shaded, the mechanical-versus-dynamical decomposition
of Section~\ref{sec:mechanical}, and the struct/align/relabel split. Not yet
run}
